% Annexe de la partie IR-CSG (réécrite le 16/11/2026 ; base 2025, taux marginaux recalculés et cotisations retirées le
% 17/11/2026 ; coupée le 17/11/2026, CONSIGNES_COUPE.md). Chiffres : sorties/*_ircsg.txt, controles_purs.txt,
% controle_hauts_purs.txt, architectures_ircsg.txt ; tableaux : tables_ircsg.py ; figures : figures_ircsg_n.py.

\subsection{Champ, données et modèle}

\paragraph{Le champ.} En 2025, l'IR, la CSG et la CRDS prélèvent 229,9~Md€ sur les revenus imposés au barème : 86,6~Md€
d'IR, 135,8~Md€ de CSG (108,5 sur les revenus d'activité, 26,7 sur les pensions, 0,6 sur le chômage) et 7,6~Md€ de CRDS.
Les revenus du capital au prélèvement forfaitaire unique gardent leur régime (section~\ref{sec:pfu}) ; on leur impute
11,5~Md€ de l'IR de l'IRCOM des revenus 2024 (12,8~\% de 90~Md€ de revenus du capital hors barème, convention), 12,4~Md€
en 2025. Ces 90~Md€ comprenant 7,6~Md€ d'assurance-vie et 1,7~Md€ de plus-values immobilières, hors PFU, l'IR imputé est
surestimé d'environ 1,4~Md€ : corrigé, la recette à remplacer passerait à 231,3~Md€ et le 1~\% paierait environ 1,1~Md€
de moins (modèle non recalé ; section~\ref{sec:ens_haut}). La CASA, le prélèvement de solidarité sur les loyers et la
CEHR restent hors champ.

\paragraph{La base 2025.} La répartition est celle des revenus de 2024 (IRCOM), chaque montant étant porté en 2025 avec
la croissance de son agrégat, la même dans toutes les tranches : IR des comptes nationaux (96,0~Md€ en 2024, 103,6~Md€
en 2025, $+$7,9~\%), CSG sur les revenus d'activité ($+$2,3~\%) et de remplacement ($+$3,3~\%) (rapport « Financement »
annexé au projet de loi d'approbation des comptes de la Sécurité sociale de 2025), CRDS ($+$2,2~\%) ; les revenus
grandissent comme l'assiette de la CSG qui les frappe. La hausse de l'IR encaissé tient en partie au jeu des acomptes et
des soldes (Insee) ; le budget de l'État la confirme (IR net de 94,6~Md€, en hausse de 7,0~Md€, selon la commission des
finances du Sénat). Si elle est passagère, la recette à remplacer est surestimée d'environ 4,5~Md€ (0,4 point de taux).
Restent en 2024 : la répartition entre foyers, la composition des revenus, la CSG des pensions de départ, les niches et
le taux marginal actuel par tranche. Le droit de l'IR est celui des revenus 2025.

\paragraph{Les données et les foyers.} Les 25 tranches de revenu fiscal de l'IRCOM sont découpées en 60 foyers
représentatifs chacune ; la composition des revenus vient des totaux par case des déclarations 2042, la CSG d'une
reconstitution calée sur les comptes de la Sécurité sociale (25,4~Md€ sur les pensions en 2024, chiffre de la CCSS non
relu à la source). L'assiette est de 1\,520~Md€ après cotisations en 2025, dont 420~Md€ de pensions. Chaque foyer est
scindé en un foyer d'actifs et un foyer de retraités de même revenu, dans la proportion des pensions de sa tranche :
27,1~millions de foyers d'actifs, 14,5~millions de foyers de retraités. Les couples partagent leurs revenus selon
l'Insee (Insee Première 1492) ; les enfants ne sont que chez les actifs (un par parent isolé, deux par couple « avec
enfants »). La charge actuelle de chaque tranche est la charge observée, répartie entre types de foyers selon des
cas-types OpenFisca sans niche ; les retraités paient ainsi 50,2~Md€, 12,0~\% de leurs pensions, et les actifs 16,3~\%
de leurs revenus (ni la DGFiP ni la DREES ne publient l'impôt des foyers de retraités). Ce sont des ordres de grandeur,
pas une microsimulation : la dispersion réelle des gagnants et des perdants est plus forte.

\paragraph{Le nouvel impôt.} Les options du modèle (\texttt{modele\_purs.impot\_adultes}) sont la franchise d'activité
(salaires, indépendants et loyers, que le modèle ne sait pas séparer), l'exonération des petites pensions (totale
jusqu'à 12\,500~€, à peu près le minimum vieillesse de 12\,523~€ en 2026, puis réduite d'un euro par euro de pension,
d'un demi-euro dans la variante « étalée »), l'abattement sur les pensions, la franchise ou le chèque par enfant, les
quotients aux plafonds actuels et la déduction de l'épargne retraite (10,8~Md€ de versements en 2025), supposée
proportionnelle aux revenus d'activité de la moitié la plus aisée, faute de répartition publiée ; elle coûte 3,0~Md€,
soit 0,29 point de taux pour tous.

\paragraph{L'efficacité.} Le gain est la baisse de la perte sèche statique sur l'activité, à court terme (méthode v2,
\texttt{revue/methode\_v2.md}). Seul le revenu d'activité des actifs réagit ; chaque euro perdu coûte tout ce que la
collectivité prélevait dessus (annexe de méthode). La participation réagit au taux de prélèvement du premier salaire,
sous deux Smic ; la réaction du revenu déclaré du dixième le plus aisé est comptée comme un effet sur la recette
(tableau~\ref{tab:ircsg_jeux}).

\begin{table}[!htbp]\centering\scriptsize\setlength{\tabcolsep}{4pt}
\caption{Les jeux d'élasticités}\label{tab:ircsg_jeux}
\begin{tabular}{@{}lrrrr@{}}
\toprule
& Bas & Central & Haut & Long terme \\
\midrule
Revenu réel, moitié du bas & 0,02 & 0,05 & 0,10 & 0,10 \\
Revenu réel, 50 à 90~\% & 0,05 & 0,10 & 0,20 & 0,20 \\
Revenu réel, 90 à 99~\% et 1~\% du haut & 0,05 & 0,10 & 0,25 & 0,25 \\
Revenu déclaré, 90 à 99~\% / 1~\% (effet sur la recette) & 0,10 / 0,15 & 0,20 / 0,30 & 0,30 / 0,45 & 0,20 / 0,30 \\
Participation, sous deux Smic & 0,10 & 0,25 & 0,50 & 0,10 \\
\bottomrule
\end{tabular}
\par\vspace{3pt}\parbox{\linewidth}{\raggedright\scriptsize\color{encre}
Appuis : Lehmann, Marical et Rioux (2013), Cabannes, Houdré et Landais (2014), Sicsic (2022), Chetty, Guren, Manoli et
Weber (2013) ; \texttt{cours/elasticites-france.tex}. Ce sont des choix, pas des estimations. Jeu de long terme :
\texttt{revue/avis\_cochrane.md}. Déplacement du revenu déclaré vers l'IS et le PFU : la moitié, à 48,6~\%
(convention).}
\end{table}

\paragraph{Les cas-types et les taux marginaux.} OpenFisca-France 176.1, législation des revenus 2025
(\texttt{openfisca\_cas\_ircsg.py}) : salariés du privé, retraités nés en 1950, aucune niche, couples mariés ou pacsés,
Smic de 21\,622~€ brut par an (sources secondaires). Les seuils de CSG des retraités y sont de 12\,817, 16\,755 et
26\,004~€ de revenu fiscal pour une part, non relus ; ceux de 2026 (13\,048, 17\,057 et 26\,472~€) ont été vérifiés sur
service-public.fr. Le taux de CSG dépendant du revenu fiscal de l'avant-dernière année, qui dépend de la CSG déductible,
on suppose la même pension chaque année et l'on retient le taux le plus élevé dont le seuil est dépassé par le revenu
fiscal calculé à ce taux. Les taux marginaux sont calculés tous les 100~€, sans lissage ; leurs pics sont de vrais sauts
du droit : seuil de CSG appliqué à toute la pension, abattement des plus de 65 ans qui baisse d'un coup de 1\,411~€ à
chacun de ses deux seuils (vers 33\,500~€ de pension, 100~€ de plus coûtent 247~€) et, dans la proposition, sortie de
l'exonération des petites pensions.

\subsection{Situation actuelle et réforme naïve : compléments}

\figbloc{ircsg_n02_ocde}{La France prélève sur le revenu un peu plus que la moyenne de l'OCDE, pour près de 60~\% par
la CSG et la CRDS}{fig:ircsg_ocde}

\figbloc{ircsg_n05_quotient}{Le quotient familial ne vaut rien sous 45\,000~€ de revenu, et jusqu'à 7\,200~€ par an
au-delà de 114\,000~€}{fig:ircsg_quotient}

\figbloc{ircsg_n07_naive_marges}{La fusion naïve garde des marches, à d'autres endroits}{fig:ircsg_naive_marges}

Appliqué à chaque adulte sur l'assiette large, le barème de l'IR rapporte 205~Md€ et le taux proportionnel tombe à
1,64~\%. Le dixième le plus aisé paie 2\,950~€ de plus par foyer (12,5~Md€), le 1~\% le plus aisé 2,8~Md€ ; un couple
avec deux enfants à un seul revenu de 50\,000~€, 2\,360~€ de plus. Le gain est de 2,6~Md€ (1,1 à 5,6), presque
entièrement sur la participation. Le passage de l'IR actuel (86,6~Md€) à ce barème se décompose ainsi : règles actuelles
sans niche, 115,1~Md€ ; assiette large, $+$47,5 ; sans quotient familial, $+$17,4 ; sans décote, $+$7,0 ; sans quotient
conjugal, $+$18,0. L'écart de 28,5~Md€ entre 115,1 et 86,6 recouvre les niches (au moins 24~Md€) et des revenus exonérés
présents dans le revenu fiscal ; il n'est pas bouclé.

\subsection{Le choix de la proposition}

La recherche porte sur la franchise d'activité (0 à 8\,000~€ par adulte), le taux d'entrée (0 à 16~\%), le taux à
20\,000~€ (8 à 26~\%) et le seuil intermédiaire (25\,000 à 90\,000~€). Le taux au seuil intermédiaire est calé sur la
recette, le seuil du taux de 50~\% sur une charge inchangée du 1~\% le plus aisé (à 0,05~Md€ près). Sur 10\,332 barèmes
essayés, 3\,743 tiennent la recette et la charge du 1~\%, et 1\,119 laissent aussi chaque décile à moins de 500~€ de sa
charge actuelle (figure~\ref{fig:ircsg_deciles}). La règle de choix, le gain sur la participation le plus élevé (à
0,02~Md€ près) puis le gain total, retient une franchise de 8\,000~€, à la borne de l'intervalle, et un taux d'entrée de
15~\%. Le barème le plus efficace a la même franchise et un taux d'entrée de 16~\% : il gagne autant (4,69~Md€ contre
4,67), mais 2,1~Md€ seulement sur la participation (tableau~\ref{tab:ircsg_design}, figure~\ref{fig:ircsg_designs}). Le
barème retenu passe par 15~\% à 20\,000~€, 28,8~\% à 30\,000~€, 30~\% à 45\,000~€, 36~\% à 100\,000~€, 41,5~\% à
150\,000~€ et 50~\% à 231\,000~€.

\figbloc{ircsg_n10_deciles}{Avec la proposition, aucun décile ne bouge de plus de 500~€}{fig:ircsg_deciles}

\begin{table}[!htbp]\centering\scriptsize\setlength{\tabcolsep}{3pt}
\caption{Franchise d'activité contre taux d'entrée : le meilleur barème pour chaque combinaison}\label{tab:ircsg_design}
\resizebox{\linewidth}{!}{%
\begin{tabular}{@{}rrrrrrrl@{}}
\toprule
Franchise & Taux & Gain & dont & Retraités & \multicolumn{2}{c}{Premier salaire (\%)} & Barème \\ \cmidrule(lr){6-7}
(€) & d'entrée & Md€ & participation & Md€ & mi-temps & Smic & \\
\midrule
0 & 0~\% & 3,26 & 0,88 & +6,9 & 4,6 & 9,2 & 22~\% à 20~k€, 22,5~\% à 30~k€, 50~\% à 212~k€ \\
0 & 2~\% & 3,49 & 0,78 & +6,7 & 5,7 & 9,7 & 21~\% à 20~k€, 21,2~\% à 30~k€, 50~\% à 194~k€ \\
0 & 4~\% & 3,69 & 0,82 & +6,6 & 6,7 & 9,8 & 19~\% à 20~k€, 24,4~\% à 45~k€, 50~\% à 201~k€ \\
0 & 6~\% & 3,89 & 0,69 & +6,4 & 7,8 & 10,3 & 18~\% à 20~k€, 21,7~\% à 37,5~k€, 50~\% à 184~k€ \\
0 & 8~\% & 4,07 & 0,76 & +6,2 & 8,8 & 10,4 & 16~\% à 20~k€, 24,5~\% à 45~k€, 50~\% à 194~k€ \\
0 & 10~\% & 4,20 & 0,96 & +6,2 & 9,5 & 10,1 & 13~\% à 20~k€, 25,2~\% à 37,5~k€, 50~\% à 219~k€ \\
2\,000 & 0~\% & 3,26 & 1,22 & +10,6 & 3,1 & 8,0 & 24~\% à 20~k€, 28,8~\% à 90~k€, 50~\% à 176~k€ \\
2\,000 & 2~\% & 3,51 & 1,28 & +10,3 & 4,0 & 8,3 & 22~\% à 20~k€, 26,3~\% à 45~k€, 50~\% à 231~k€ \\
2\,000 & 4~\% & 3,74 & 1,16 & +10,1 & 5,0 & 8,9 & 21~\% à 20~k€, 32,3~\% à 90~k€, 50~\% à 212~k€ \\
2\,000 & 6~\% & 3,96 & 1,20 & +9,9 & 5,9 & 9,1 & 19~\% à 20~k€, 24,6~\% à 37,5~k€, 50~\% à 212~k€ \\
2\,000 & 8~\% & 4,17 & 1,23 & +9,8 & 6,8 & 9,4 & 17~\% à 20~k€, 25,3~\% à 37,5~k€, 50~\% à 219~k€ \\
2\,000 & 10~\% & 4,36 & 1,26 & +9,6 & 7,6 & 9,6 & 15~\% à 20~k€, 26,0~\% à 37,5~k€, 50~\% à 227~k€ \\
2\,000 & 12~\% & 4,56 & 1,13 & +9,3 & 8,6 & 10,2 & 14~\% à 20~k€, 24,7~\% à 37,5~k€, 50~\% à 205~k€ \\
4\,000 & 0~\% & 3,30 & 1,62 & +14,4 & 1,8 & 6,8 & 26~\% à 20~k€, 28,8~\% à 90~k€, 50~\% à 182~k€ \\
4\,000 & 2~\% & 3,52 & 1,50 & +14,2 & 2,6 & 7,4 & 25~\% à 20~k€, 27,4~\% à 90~k€, 50~\% à 156~k€ \\
4\,000 & 4~\% & 3,78 & 1,54 & +13,9 & 3,4 & 7,7 & 23~\% à 20~k€, 32,0~\% à 90~k€, 50~\% à 212~k€ \\
4\,000 & 6~\% & 4,01 & 1,67 & +13,8 & 4,0 & 7,8 & 20~\% à 20~k€, 25,9~\% à 30~k€, 50~\% à 247~k€ \\
4\,000 & 8~\% & 4,23 & 1,57 & +13,4 & 4,8 & 8,4 & 19~\% à 20~k€, 26,4~\% à 37,5~k€, 50~\% à 231~k€ \\
4\,000 & 10~\% & 4,44 & 1,58 & +13,2 & 5,6 & 8,8 & 17~\% à 20~k€, 27,0~\% à 37,5~k€, 50~\% à 239~k€ \\
4\,000 & 12~\% & 4,64 & 1,45 & +12,9 & 6,4 & 9,4 & 16~\% à 20~k€, 25,7~\% à 37,5~k€, 50~\% à 215~k€ \\
4\,000 & 14~\% & 4,83 & 1,45 & +12,7 & 7,2 & 9,8 & 14~\% à 20~k€, 26,4~\% à 37,5~k€, 50~\% à 223~k€ \\
6\,000 & 6~\% & 4,05 & 1,93 & +17,6 & 2,4 & 6,9 & 23~\% à 20~k€, 25,2~\% à 25~k€, 50~\% à 239~k€ \\
6\,000 & 8~\% & 4,28 & 1,92 & +17,3 & 2,9 & 7,4 & 21~\% à 20~k€, 26,1~\% à 30~k€, 50~\% à 239~k€ \\
6\,000 & 10~\% & 4,50 & 1,91 & +17,1 & 3,5 & 7,8 & 19~\% à 20~k€, 26,2~\% à 30~k€, 50~\% à 239~k€ \\
6\,000 & 12~\% & 4,71 & 2,00 & +17,0 & 4,0 & 8,0 & 16~\% à 20~k€, 27,8~\% à 30~k€, 50~\% à 274~k€ \\
6\,000 & 14~\% & 4,89 & 1,87 & +16,6 & 4,5 & 8,7 & 15~\% à 20~k€, 26,4~\% à 30~k€, 50~\% à 239~k€ \\
7\,000 & 6~\% & 4,08 & 2,03 & +19,6 & 1,6 & 6,5 & 25~\% à 20~k€, 31,3~\% à 90~k€, 50~\% à 200~k€ \\
7\,000 & 8~\% & 4,28 & 2,01 & +19,3 & 2,0 & 6,9 & 23~\% à 20~k€, 25,5~\% à 30~k€, 50~\% à 227~k€ \\
7\,000 & 10~\% & 4,53 & 2,09 & +19,1 & 2,4 & 7,2 & 20~\% à 20~k€, 27,1~\% à 30~k€, 50~\% à 256~k€ \\
7\,000 & 12~\% & 4,73 & 2,07 & +18,9 & 2,8 & 7,7 & 18~\% à 20~k€, 27,2~\% à 30~k€, 50~\% à 256~k€ \\
7\,000 & 14~\% & 4,89 & 2,04 & +18,6 & 3,2 & 8,1 & 16~\% à 20~k€, 27,2~\% à 30~k€, 50~\% à 247~k€ \\
7\,000 & 16~\% & 5,02 & 1,80 & +18,2 & 3,7 & 8,9 & 16~\% à 20~k€, 24,3~\% à 30~k€, 50~\% à 201~k€ \\
8\,000 & 8~\% & 4,31 & 2,21 & +21,4 & 1,2 & 6,3 & 24~\% à 20~k€, 25,8~\% à 25~k€, 50~\% à 239~k€ \\
8\,000 & 10~\% & 4,51 & 2,19 & +21,1 & 1,4 & 6,8 & 22~\% à 20~k€, 26,5~\% à 30~k€, 50~\% à 239~k€ \\
8\,000 & 12~\% & 4,72 & 2,23 & +21,0 & 1,6 & 7,1 & 19~\% à 20~k€, 26,4~\% à 25~k€, 50~\% à 247~k€ \\
8\,000 & 14~\% & 4,91 & 2,27 & +20,8 & 1,9 & 7,4 & 16~\% à 20~k€, 27,4~\% à 25~k€, 50~\% à 265~k€ \\
8\,000 & 16~\% & 5,04 & 2,06 & +20,3 & 2,2 & 8,2 & 16~\% à 20~k€, 24,7~\% à 25~k€, 50~\% à 215~k€ \\
\bottomrule
\end{tabular}}
\par\vspace{3pt}\parbox{\linewidth}{\raggedright\scriptsize\color{encre}
Pour chaque franchise d'activité (par adulte) et chaque taux d'entrée (taux sur le premier euro imposé après la franchise), le barème le plus efficace parmi ceux qui respectent la recette, les déciles à $\pm$500~€ et la neutralité du 1~\%. Sont aussi fixées : l'exonération des petites pensions, la franchise de 6\,000~€ par enfant, la déduction de l'épargne retraite. Premier salaire : taux de prélèvement IR-CSG d'une personne seule qui prend un emploi au Smic, à mi-temps ou à temps plein (aujourd'hui : 9,5~\%). Taux d'entrée pairs seulement ; 10332 barèmes essayés, 1047 admissibles. La proposition retient la combinaison qui maximise la participation (franchise de 8\,000~€, taux d'entrée de 14~\%).}
\end{table}


\figbloc{ircsg_n17_designs}{Plus la franchise est forte, plus la participation progresse, et le gain total avec
elle}{fig:ircsg_designs}

\subsection{Les variantes}

Le tableau~\ref{tab:ircsg_dec} donne la variation par décile de toutes les variantes, le tableau~\ref{tab:ircsg_arbcas}
et la figure~\ref{fig:ircsg_cas} leurs effets sur des ménages types, la figure~\ref{fig:ircsg_arbitrages} le gain de
chaque arbitrage ; le tableau~\ref{tab:ircsg_arch} rappelle les architectures de la méthode v2, où chaque dispositif est
ajouté seul au barème neutre et financé par des points de taux pour tous. Un chèque par enfant égal au coût actuel du
quotient familial (765~€) libérerait 0,7 point de taux pour tous.

\begin{table}[!htbp]\centering\scriptsize\setlength{\tabcolsep}{3pt}
\caption{Variation de l'impôt par décile de revenu fiscal, toutes variantes (€ par foyer et par an)}\label{tab:ircsg_dec}
\resizebox{\linewidth}{!}{%
\begin{tabular}{@{}lrrrrrrrrrrrr@{}}
\toprule
& D1 & D2 & D3 & D4 & D5 & D6 & D7 & D8 & D9 & D10 & 1~\% (Md€) & Retraités (Md€) \\
\midrule
Réforme naïve (quotient, 3,99~\%) & $-$11 & $-$37 & +162 & +105 & $-$543 & $-$551 & $-$570 & +95 & $-$524 & +1\,841 & +2,6 & +8,2 \\
Réforme naïve, variante individuelle (1,64~\%) & $-$57 & $-$217 & $-$171 & $-$256 & $-$906 & $-$938 & $-$794 & +36 & +277 & +2\,953 & +2,8 & +2,7 \\
\textbf{Proposition} & $-$84 & $-$279 & +80 & +251 & $-$215 & $-$155 & $-$320 & +354 & $-$71 & +436 & 0,0 & +21,5 \\
Variante : franchise de 4\,000~€ & $-$85 & $-$149 & +147 & +225 & $-$286 & $-$129 & $-$237 & +304 & $-$158 & +367 & 0,0 & +14,3 \\
Variante : sans franchise & +2 & +19 & +243 & +318 & $-$225 & $-$202 & $-$414 & +113 & $-$259 & +399 & 0,0 & +6,9 \\
Quotient conjugal et familial & $-$84 & $-$285 & $-$18 & +119 & $-$276 & $-$137 & $-$149 & +468 & $-$99 & +460 & 0,0 & +18,9 \\
Abattement de 10~\% sur les pensions & $-$83 & $-$277 & +42 & +273 & $-$200 & $-$138 & $-$287 & +314 & $-$94 & +448 & 0,0 & +14,8 \\
Chèque par enfant (même coût) & $-$168 & $-$354 & +16 & +197 & $-$226 & $-$99 & $-$227 & +392 & +11 & +455 & 0,0 & +21,6 \\
Le 1~\% peut payer moins & $-$84 & $-$278 & +105 & +335 & $-$86 & $-$11 & $-$197 & +493 & +162 & $-$423 & $-$3,8 & +22,5 \\
Barème neutre (méthode v2) & 0 & +83 & +265 & +165 & $-$365 & $-$221 & $-$324 & +130 & $-$27 & +290 & $-$0,5 & +5,0 \\
Ancienne proposition (1~\% en baisse) & $-$85 & $-$145 & +190 & +342 & $-$150 & $-$12 & $-$121 & +450 & 0 & $-$453 & $-$2,5 & +13,9 \\
\bottomrule
\end{tabular}}
\par\vspace{3pt}\parbox{\linewidth}{\raggedright\scriptsize\color{encre}
Déciles de foyers fiscaux classés par revenu fiscal de référence (IRCOM, revenus 2024 ; montants de 2025). Moyennes de foyers représentatifs : elles cachent les écarts à l'intérieur de chaque décile (sous-partie~\ref{sec:ircsg}.4).}
\end{table}


\begin{table}[!htbp]\centering\scriptsize\setlength{\tabcolsep}{3pt}
\caption{Les arbitrages pour des ménages types (variation de l'IR + CSG + CRDS, € par an)}\label{tab:ircsg_arbcas}
\resizebox{\linewidth}{!}{%
\begin{tabular}{@{}lrrrrr@{}}
\toprule
& Proposition & Quotient & Abatt. 10~\% & Chèque & Franchise 4\,000~€ \\
\midrule
Jeune salarié au Smic & $-$447 & $-$635 & $-$259 & $-$359 & $-$243 \\
Cadre seul, 100 000 € brut & $-$4\,015 & $-$2\,733 & $-$3\,682 & $-$4\,962 & $-$3\,894 \\
Parent isolé, deux enfants, au Smic & $-$2\,027 & $-$688 & $-$2\,027 & $-$2\,740 & $-$1\,683 \\
Couple sans enfant à un seul revenu, 80 000 € & +3\,498 & +65 & +3\,814 & +2\,873 & +3\,677 \\
Couple avec deux enfants à un seul revenu, 50 000 € & $-$862 & $-$77 & $-$539 & +92 & $-$518 \\
Couple avec trois enfants, 60 000 € brut (60/40) & $-$2\,922 & $-$993 & $-$2\,591 & $-$3\,078 & $-$2\,452 \\
Couple avec trois enfants à un seul revenu, 60 000 € & $-$1\,003 & +102 & $-$682 & +315 & $-$670 \\
Retraité seul, pension de 16 000 € & +1\,005 & +877 & +871 & +1\,041 & +840 \\
Retraité seul, pension de 20 000 € & +1\,353 & +1\,109 & +1\,262 & +1\,521 & +940 \\
Retraité seul, pension de 25 000 € & +1\,815 & +1\,577 & +1\,519 & +2\,118 & +1\,092 \\
Couple de retraités, 45 000 € de pensions (60/40) & +1\,330 & +196 & +809 & +1\,667 & +309 \\
\bottomrule
\end{tabular}}
\par\vspace{3pt}\parbox{\linewidth}{\raggedright\scriptsize\color{encre}
Mêmes cas-types que le tableau~\ref{tab:ircsg_cas}. Quotient : imposition commune des couples, quotient conjugal et familial aux plafonds actuels à la place de la franchise par enfant ; abattement de 10~\% sur les pensions en plus de l'exonération des petites pensions ; chèque de 1\,191~€ par enfant, de même coût que la franchise de 6\,000~€ par enfant ; franchise d'activité de 4\,000~€ par adulte au lieu de 8\,000~€. Chaque variante a son propre barème, recherché sous les mêmes contraintes.}
\end{table}


\figbloc{ircsg_n15_cas_types}{Familles et salariés gagnent, retraités et couples à un seul revenu perdent}{fig:ircsg_cas}

\figbloc{ircsg_n14_arbitrages}{Chaque arbitrage déplace le gain de moins d'un milliard, mais la charge des retraités de
plusieurs}{fig:ircsg_arbitrages}

\figbloc{ircsg_n18_variantes_marges}{Étaler la sortie de l'exonération abaisse le taux des petites pensions et le
relève un peu plus haut}{fig:ircsg_variantes_marges}

\begin{table}[!htbp]\centering\scriptsize\setlength{\tabcolsep}{3pt}
\caption{Architectures comparées, à recette égale (220~Md€)}\label{tab:ircsg_arch}
\resizebox{\linewidth}{!}{%
\begin{tabular}{@{}lrrrrrrr@{}}
\toprule
& Points & Gain d'efficacité & Part du & Perdants & Retraités & Couple, & 10~\% \\
& pour tous & Md€ par an & gain & $>$ 500~€ & Md€ & 2 enfants & du haut \\
& & (fourchette) & atteignable & & & € par foyer & Md€ \\
\midrule
Barème neutre & 0,0 & 4,0 (1,9 à 9,2) & 52~\% & 25~\% & +7,0 & +570 & +2,5 \\
\multicolumn{8}{@{}l}{\emph{Un dispositif ajouté au barème neutre}} \\
\quad Franchise d'activité 4~000~€ & +1,8 & 4,6 (2,1 à 10,2) & 59~\% & 30~\% & +14,0 & +217 & +3,6 \\
\quad Franchise d'activité 8~000~€ & +3,6 & 5,0 (2,2 à 11,0) & 64~\% & 33~\% & +21,4 & $-$172 & +4,7 \\
\quad Franchise de 8~000~€ sur tous les revenus & +5,8 & 3,3 (1,4 à 7,4) & 43~\% & 21~\% & +5,4 & +775 & +9,6 \\
\quad Franchise de 6~000~€ par enfant & +1,5 & 3,6 (1,7 à 8,4) & 47~\% & 27~\% & +13,0 & $-$983 & +0,6 \\
\quad Chèque de 742~€ par enfant & +0,9 & 3,1 (1,5 à 7,3) & 40~\% & 26~\% & +10,7 & $-$361 & +3,1 \\
\quad Quotient conjugal et familial & +2,0 & 3,9 (1,9 à 9,0) & 50~\% & 28~\% & +12,1 & $-$965 & $-$3,0 \\
\quad Abattement de 10~\% sur les pensions & +0,6 & 3,4 (1,6 à 8,0) & 44~\% & 24~\% & +0,9 & +923 & +2,7 \\
\quad Petites pensions exonérées & +0,4 & 3,7 (1,7 à 8,5) & 47~\% & 16~\% & +3,2 & +786 & +4,0 \\
\midrule
\textbf{Proposition}, franchise d'activité 4~000~€ & --- & 4,9 (2,3 à 11,1) & 63~\% & 29~\% & +15,2 & $-$1\,195 & $-$1,7 \\
Proposition, franchise d'activité 8~000~€ & --- & 5,2 (2,4 à 11,8) & 67~\% & 26~\% & +22,8 & $-$1\,629 & $-$1,5 \\
\midrule
Optimum d'efficacité, sans franchise & --- & 7,4 (3,7 à 17,4) & 95~\% & 58~\% & +13,7 & $-$519 & $-$25,7 \\
Gain atteignable (optimum, franchise 4~000~€) & --- & 7,8 (3,8 à 18,2) & 100~\% & 49~\% & +20,3 & $-$869 & $-$25,1 \\
\bottomrule
\end{tabular}}
\par\vspace{3pt}\parbox{\linewidth}{\raggedright\scriptsize\color{encre}
Gain d'efficacité : baisse de la perte sèche sur l'activité réelle par rapport à aujourd'hui, hypothèses centrales (basses et hautes entre parenthèses) ; la réaction du revenu déclaré n'y est pas comptée. Points pour tous : hausse des taux du barème neutre qui finance le dispositif. Proposition : franchise d'activité, 6~000~€ par enfant, pensions exonérées jusqu'à 12~500~€ puis de moins en moins jusqu'à 25~000~€ ; barème le plus efficace parmi ceux qui ne changent la charge moyenne d'aucun décile de plus de 500~€ et n'alourdissent pas celle du 1~\% le plus aisé. Retraités, couples, 10~\% du haut : variation de leur impôt (positif : ils paient plus).}
\end{table}


\subsection{Sensibilités, haut de la distribution et transition}

\paragraph{Ce qui fait varier le gain.} La perception de la CSG domine toutes les autres sensibilités
(tableau~\ref{tab:ircsg_sens}) : la fusion n'est efficace que si la CSG d'aujourd'hui fait déjà réagir les salariés.

\begin{table}[!htbp]\centering\scriptsize\setlength{\tabcolsep}{4pt}
\caption{Ce qui fait varier le gain de la proposition et la neutralité du 1~\%}\label{tab:ircsg_sens}
\resizebox{\linewidth}{!}{%
\begin{tabular}{@{}lrl@{}}
\toprule
Hypothèse & Gain (Md€) & Autre effet \\
\midrule
Modèle, hypothèses centrales & 4,91 & participation 2,27~Md€ \\
Élasticités basses / hautes & 2,25 / 10,94 & \\
Élasticités de long terme (intensif 0,10 à 0,25 ; participation 0,10) & 7,30 & participation 0,91~Md€ \\
CSG perçue à moitié aujourd'hui & 1,29 & intensif $-$0,98~Md€ \\
CSG non perçue aujourd'hui & $-$2,09 & intensif $-$4,37~Md€ \\
Régime permanent : impôt futur sur les pensions vu à moitié par les actifs & 4,10 & \\
Régime permanent : vu en entier & 3,25 & coin $+$1,40~point \\
Répartition des revenus au PFU logés en haut (convention « haut ») & 4,5 & 1~\% $-$2,8~Md€ ; neutre si 50~\% à 110~k€ \\
Répartition des revenus au PFU au prorata & 5,3 & 1~\% +3,7~Md€ ; neutre si 50~\% à 858~k€ \\
\bottomrule
\end{tabular}}
\par\vspace{3pt}\parbox{\linewidth}{\raggedright\scriptsize\color{encre}
Gain : baisse de la perte sèche sur l'activité (jeu central sauf mention). CSG non perçue : le revenu d'activité ne réagit aujourd'hui qu'à l'IR (comme aux cotisations chez Lehmann, Marical et Rioux, 2013) et réagit après la fusion à tout l'impôt fusionné ; la marge de participation est gardée. Régime permanent : la hausse de l'impôt sur les pensions réduit la valeur des droits acquis par les cotisations contributives (27,4~\% du salaire brut) ; borne haute, le lien entre cotisation et pension étant faible au bas de l'échelle et nul au-delà de 8 plafonds. Convention centrale de répartition des revenus au PFU : « mixte » (controle\_hauts\_purs.py).}
\end{table}


\paragraph{Le haut de la distribution.} La DGFiP ne publie pas, par tranche, la part des revenus imposés au PFU. Dans la
convention centrale (« mixte »), la charge du 1~\% le plus aisé ne bouge pas : la baisse des hauts salaires y est
compensée par une hausse de 3 à 16 points pour les revenus du barème au-delà de 1~M€, auxquels cette convention prête un
taux actuel de 33 à 39~\%. Dans la convention « haut », plus proche du barème légal (44 à 51~\%) et que soutient un
calage sur l'IR observé (\texttt{revue/hypotheses\_ircsg.md}, H9), le 1~\% paierait 3,0~Md€ de moins, les foyers de 1 à
9~M€ de 2 à 8 points de moins ; garder sa charge demanderait le taux de 50~\% dès 91\,000~€ environ, mais le dixième le
plus aisé paierait alors 1\,300~€ de plus en moyenne. Au prorata, le 1~\% paierait 3,8~Md€ de plus
(\texttt{sorties/controle\_hauts\_purs.txt}). La disparition de la CEHR (3 à 4 points au très haut) sur les revenus du
barème, qui rend 0,3 à 0,7~Md€ au 1~\%, n'est pas comptée ici mais dans le bilan du haut (section~\ref{sec:ens_haut}).

\paragraph{La transition.} Le tableau~\ref{tab:ircsg_trans} donne la hausse annuelle des principaux perdants selon la
durée de la convergence.

\begin{table}[!htbp]\centering\scriptsize\setlength{\tabcolsep}{4pt}
\caption{Hausse annuelle des principaux perdants selon la durée de la transition (€ par an)}\label{tab:ircsg_trans}
\begin{tabular}{@{}lrrrr@{}}
\toprule
& Hausse finale & \multicolumn{3}{c}{Hausse chaque année, transition sur} \\ \cmidrule(l){3-5}
& & 3 ans & 5 ans & 10 ans \\
\midrule
Retraités, médiane & +1\,392 & 464 & 278 & 139 \\
Retraités, 90\textsuperscript{e} centile & +3\,423 & 1\,141 & 685 & 342 \\
Couples de retraités, 90\textsuperscript{e} centile & +5\,659 & 1\,886 & 1\,132 & 566 \\
Retraités, 99\textsuperscript{e} centile & +7\,423 & 2\,474 & 1\,485 & 742 \\
Tous les foyers, 90\textsuperscript{e} centile & +2\,245 & 748 & 449 & 224 \\
Retraité seul, pension de 16 000 € & +1\,050 & 350 & 210 & 105 \\
Retraité seul, pension de 20 000 € & +1\,390 & 463 & 278 & 139 \\
Retraité seul, pension de 25 000 € & +1\,679 & 560 & 336 & 168 \\
Couple de retraités, 45 000 € de pensions (60/40) & +1\,834 & 611 & 367 & 183 \\
Couple sans enfant à un seul revenu, 80 000 € & +3\,715 & 1\,238 & 743 & 372 \\
\bottomrule
\end{tabular}
\par\vspace{3pt}\parbox{\linewidth}{\raggedright\scriptsize\color{encre}
Convergence linéaire : chaque année, chaque foyer paie une part croissante du nouvel impôt et décroissante de l'ancien, ce qui tient la recette chaque année. Centiles : foyers représentatifs du modèle (la dispersion réelle des hausses est plus forte ; le 99\textsuperscript{e} centile en est un minorant). Cas-types : tableau~\ref{tab:ircsg_cas}.}
\end{table}


\paragraph{La valeur actualisée.} La hausse annuelle des retraités est multipliée par un facteur d'annuité : table de
Gompertz ($\mu(x) = A\,e^{0,10x}$) calée sur une espérance de vie à 65 ans de 21,8 ans (Insee 2025 provisoire) ;
cohortes de même taille dès 62 ans (âge moyen 76 ans, hypothèse) ; taux réel de 1 à 3~\%. Le facteur moyen va de 10,8 à
12,9 (8,6 à 10,7 avec une montée en charge sur cinq ans). Au taux de 2~\%, un euro de hausse annuelle vaut 17~€ à 65
ans, 12~€ à 75 ans et 7~€ à 85 ans (\texttt{sorties/transition\_ircsg.txt}).

\paragraph{Le gain permanent.} En régime permanent, la hausse de l'impôt sur les pensions (5,1 points de taux moyen)
réduit la valeur des droits à pension qu'ouvre chaque euro de salaire, soit environ 27 centimes par euro brut sous le
plafond de la Sécurité sociale (hypothèse) : elle alourdit d'environ 1,4 point le prélèvement implicite sur le salaire
des actifs. Selon que ceux-ci le perçoivent à moitié ou en entier, le gain permanent est de 3,8 ou de 3,0~Md€ au lieu de
4,7. C'est une borne : le lien entre un euro de salaire et la pension est faible au bas de l'échelle (minimum
contributif) et nul au-delà de huit fois le plafond.

\subsection{Hypothèses fortes et chiffres non vérifiés}

\begin{table}[!htbp]\centering\scriptsize\setlength{\tabcolsep}{3pt}
\caption{Hypothèses fortes et chiffres non vérifiés}\label{tab:ircsg_hyp}
\begin{tabular}{@{}p{0.43\linewidth}p{0.2\linewidth}p{0.33\linewidth}@{}}
\toprule
Hypothèse & Où elle joue & Ce qu'elle change si elle est fausse \\
\midrule
Perception de la CSG : la CSG et la CRDS d'aujourd'hui font réagir le revenu d'activité autant que l'IR. Aucune étude ne le mesure. & Tout le gain sur la marge intensive & Décisif. CSG non perçue : gain $-$2,5~Md€ au lieu de 4,7 ; perçue à moitié : 1,0. \\[2pt]
Élasticités : revenu réel 0,05 (moitié du bas) et 0,10 (au-dessus) ; participation 0,25 sous 2 Smic ; ce sont des choix, pas des estimations françaises. & Tous les gains & Fourchette 2,1 à 10,3~Md€ ; long terme 6,7~Md€. \\[2pt]
Les pensions, le chômage et l'épargne ne réagissent pas (âge de départ, surcote, cumul emploi-retraite, sortie en capital). & Mécanisme du financement par les pensions & Non chiffré ; sens ambigu (une pension nette plus faible pousse à partir plus tard ; un taux marginal élevé dans la zone de sortie de l'exonération pousse en sens inverse). \\[2pt]
Régime permanent : les actifs ne voient pas que leurs futures pensions seront plus taxées. & Part « permanente » du gain & Vu en entier : gain permanent 3,0~Md€ au lieu de 4,7. \\[2pt]
Répartition des revenus au PFU entre tranches (non publiée) ; convention « mixte ». & Charge du 1~\%, forme du haut du barème & Convention « haut » (que soutient un calage sur l'IR observé) : 1~\% $-$3,0~Md€ ; charge inchangée si 50~\% dès 91~k€. \\[2pt]
Foyers purs : pas de foyer mixte (couple dont un seul est retraité, retraité bailleur) ; loyers tous chez les actifs. & Coût de la franchise, charge des retraités & Franchise sous-estimée et hausse des retraités surestimée, jusqu'à 3~Md€ environ (borne grossière). \\[2pt]
Partage des revenus dans les couples (Insee Première 1492, quatre classes ; le même pour les retraités). & Coût de l'impôt individuel, charge des couples de retraités & Couples de retraités à 62/38 : charge des retraités plus faible de 2~Md€ environ (registre, H11, ancienne proposition). \\[2pt]
Enfants : un par parent isolé, deux par couple « avec enfants », aucun chez les retraités et les concubins. & Coût de la franchise par enfant, familles nombreuses & Familles de trois enfants absentes des agrégats (cas-types seulement). \\[2pt]
Charge actuelle par type de foyer répartie selon des cas-types sans niche. & Qui paie entre types de foyers & Charge actuelle des retraités (50,2~Md€) non observable ; niches réparties au prorata. \\[2pt]
Base 2025 : l'IR de 2025 est l'impôt encaissé dans l'année (comptes nationaux), dont la hausse (+7,9~\%) tient en partie au jeu des acomptes et des soldes du prélèvement à la source ; chaque agrégat de 2024 est porté en 2025 avec la même croissance dans toutes les tranches. & Recette à remplacer, charge actuelle par tranche & Si cette hausse est passagère, la recette est surestimée d'environ 4,5~Md€ (IR porté en 2025 comme les revenus), soit environ 0,4 point de taux. \\[2pt]
Épargne retraite : déduction de 10,6~Md€ (cases 2042) supposée proportionnelle aux revenus d'activité de la moitié haute. & Barème, haut de la distribution & Coût de la déduction 3,0~Md€ (0,29~point) ; répartition non vérifiée. \\[2pt]
Cas-types : législation des revenus 2025 d'OpenFisca, Smic de 21\,622~€ (sources secondaires) ; retraités en régime stable (taux de CSG le plus élevé dont le seuil est dépassé par le revenu fiscal calculé à ce taux). & Tous les cas-types, courbes de taux marginaux & Dans une bande de quelques centaines d'euros autour de chaque seuil de CSG, aucun taux n'est cohérent avec son propre revenu fiscal ; le taux inférieur est retenu, ce qui déplace le saut de moins de 1\,000~€. \\[2pt]
Valeur actualisée : table de Gompertz calée sur l'espérance de vie à 65 ans (21,8 ans) ; cohortes de retraités de même taille. & Valeur actualisée de la hausse des retraités & Ordre de grandeur ; 1 à 3~\% réel donne la fourchette. \\[2pt]
Chiffres non vérifiés sur une source primaire : CSG des pensions de 2024 (CCSS) ; croissance de la CSG par assiette en 2025 (régimes de base et FSV, appliquée à tous les affectataires) ; paramètres 2025 d'OpenFisca (seuils de CSG des pensions, abattement des plus de 65 ans) ; Smic 2025 ; part des familles au plafond du quotient familial ; partage des niches par tranche ; l'IRCOM inclut l'IR au PFU et exclut la CEHR. & Calage, cas-types, récit & Registre revue/hypotheses\_ircsg.md, section 4. \\[2pt]
\bottomrule
\end{tabular}
\par\vspace{3pt}\parbox{\linewidth}{\raggedright\scriptsize\color{encre}
Source : registre revue/hypotheses\_ircsg.md (hypothèses H1 à H20), chiffres recalculés pour la proposition par transition\_ircsg.py, controle\_hauts\_purs.py et proposition\_ircsg.py.}
\end{table}


\subsection{Programmes}

{\sloppy Dans \texttt{synthese/}, depuis \texttt{fiscalite/} : \texttt{python3.13 openfisca\_cas\_ircsg.py}, puis
\texttt{situation\_ircsg.py}, \texttt{architectures\_ircsg.py}, \texttt{naive\_ircsg.py},
\texttt{proposition\_ircsg.py}, \texttt{arbitrages\_prop\_ircsg.py}, \texttt{transition\_ircsg.py},
\texttt{controle\_hauts\_purs.py}, \texttt{tables\_ircsg.py} et \texttt{figures\_ircsg\_n.py} (deux fois, la seconde
avec \texttt{FIG\_NU=1}). Modèle : \texttt{modele\_purs.py} (\texttt{BASE\_IRCSG=2024} rend l'ancienne base) ;
outils : \texttt{outils\_ircsg.py} ; vérifications : \texttt{revue/calculs\_verif/ircsg/}.\par}
